% from ega2-1-fr.tex
\oldpage[II]{108}

rapport à $U\times_SP$ du préschéma $Z$
(\hyperref[I.9.5.8-fr]{I, 9.5.8}). Or l'immersion $\psi$ se factorise en
\[
U\xrightarrow{\Gamma_\varphi}U\times_SP\xrightarrow{j\times1}X\times_SP,
\]
et comme $P$ est séparé sur $S$, le morphisme graphe $\Gamma_\varphi$ est une
immersion fermée (\hyperref[I.5.4.3-fr]{I, 5.4.3}), donc $Z$ est un
sous-préschéma fermé de $U\times_SP$, d'où on conclut que $U'=Z$. Comme
d'ailleurs $\psi$ est une immersion, la restriction de $f$ à $U'$ est un
isomorphisme sur $U$, réciproque de $\psi'$~; enfin, par définition de $X'$,
$U'$ est dense dans $X'$.

D) Prouvons maintenant que $g$ est une immersion, ce qui établira que $X'$ est
quasi-projectif sur $S$, puisque $P$ est projectif sur $S$. Posons
\begin{align*}
V_k&=\varphi_k(U_k) &&\text{(ouvert de $P_k$)},\\
W_k&=p_k^{-1}(V_k) &&\text{(ouvert de $P$)},\\
U'_k&=f^{-1}(U_k),\qquad U''_k=g^{-1}(W_k)
&&\text{(ouverts de $X'$)}.
\end{align*}
Il est clair que les $U'_k$ forment un recouvrement ouvert de $X'$~; nous
allons d'abord voir qu'il en est de même des $U''_k$, en prouvant que
$U'_k\subset U''_k$. Il suffira pour cela de montrer que le diagramme
\[
\tag{5.6.1.5}
\label{II.5.6.1.5-fr}
\xymatrix{
U'_k\ar[r]^{g|U'_k}\ar[d]_{f|U'_k} & P\ar[d]^{p_k}\\
U_k\ar[r]_{\varphi_k} & P_k
}
\]
est commutatif. Or le préschéma $U'_k=h^{-1}(U_k\times_SP)$ est induit par
$X'$ sur l'ouvert $U'_k$, et est donc l'adhérence de $Z=U'\subset U'_k$ par
rapport à $U'_k$ (\hyperref[I.9.5.8-fr]{I, 9.5.8}). Pour montrer la
commutativité du diagramme (\hyperref[II.5.6.1.5-fr]{5.6.1.5}), il suffit
donc, puisque $P_k$ est un $S$-schéma, de montrer qu'en composant ce diagramme
avec l'injection canonique $U'\to U'_k$ (ou, ce qui revient au même en vertu de
l'isomorphisme de $U'$ et $U$, avec $\psi$), on obtient un diagramme commutatif
(\hyperref[I.9.5.6-fr]{I, 9.5.6}). Mais par définition le diagramme que l'on
obtient alors n'est autre que (\hyperref[II.5.6.1.1-fr]{5.6.1.1}), d'où notre
assertion.

Les $W_k$ forment donc un recouvrement ouvert de $g(X')$~; pour prouver que
$g$ est une immersion, il suffit de voir que chacune des restrictions
$g|U''_k$ est une immersion dans $W_k$
(\hyperref[I.4.2.4-fr]{I, 4.2.4}). Pour cela, considérons le morphisme
$u_k:W_k\xrightarrow{p_k}V_k\xrightarrow{\varphi_k^{-1}}U_k\to X$~; comme
$X$ est séparé sur $S$, le morphisme graphe
$\Gamma_{u_k}:W_k\to X\times_SW_k$ est une immersion fermée
(\hyperref[I.5.4.3-fr]{I, 5.4.3}), donc le graphe
$T_k=\Gamma_{u_k}(W_k)$ est un sous-préschéma fermé de $X\times_SW_k$~; si
nous montrons que ce sous-préschéma majore $U'$, il majorera aussi le
sous-préschéma induit par $X'$ sur l'ouvert $X''_k$ de $X'$, par
(\hyperref[I.9.5.8-fr]{I, 9.5.8}). Comme la restriction de $q_2$ à $T_k$ est
un isomorphisme sur $W_k$, la restriction de $g$ à $X''_k$ sera une immersion
dans $W_k$, et notre assertion sera démontrée. Soit $v_k$ l'injection
canonique $U'\to X\times_SW_k$~; nous devons montrer qu'il existe un morphisme
$w_k:U'\to W_k$ tel que $v_k=\Gamma_{u_k}\circ w_k$. Par définition du
produit, il suffit de prouver que
$q_1\circ v_k=u_k\circ q_2\circ v_k$
(\hyperref[I.3.2.1-fr]{I, 3.2.1}), ou, en composant à droite

